\section{制度可以废改，政治记忆并不随之消失}

1085年神宗去世，哲宗即位，高太后临朝。年幼皇帝与高度争议的新法体系，使摄政并非简单的“复旧”，而是必须逐项判断哪些法令停止、替换或保留。1086年元祐政府废罢、调整青苗、免役等多项新法。《宋会要辑稿·食货》的废改条与《长编》可追踪中枢决定，但不同制度的废止时间、地方保留与替代措施并不同步。把一纸诏令当作全国同时翻转，恰好会重复新旧两派各自最方便的叙述。

1093年高太后去世，哲宗亲政；1094年绍圣政府复行多项熙宁新法，并贬逐元祐官员。继承、摄政终结和人事转向的事实，可由《宋史》本纪、传记和会要相互核对；将其全归于皇帝个人意志，则忽略宰执、言官和外放官员形成的政策网络。制度在这里也成为官员履历的一部分：拥护或反对某项法，逐渐会被记录为可供任用、贬谪和结交的身份。

1097年继续编列、处分元祐异议官员，尤其显示这种转换如何把政策分歧变成行政分类。党人名单的版本、人数和名单中人的实际主张并不完全一致；《宋史》所用的道德化传名更不能直接等同于每位官员的政治生涯。可确定的是，名册、贬所和人事禁限使异议具有了制度化的地理与记忆。北宋中枢并非从此不能运转，却愈来愈难在政策分歧之外保留稳定的合作空间。
